\chapter{Useful Matrix Identities}

\vfill
\localtableofcontents
\vfill
\newpage

\section{Gaussian Conditioning Formula}

\begin{proposition}[Gaussian conditioning]\label{th:Gaussian_conditioning}
	For $n,p\in\mathbb{N}$, let $A\in\mathcal{M}_{n,n}(\mathbb{R})$, $B\in\mathcal{M}_{n,p}(\mathbb{R})$, and $D\in\mathcal{M}_{p,p}(\mathbb{R})$. Assuming $A\succ0$ and $D\succ0$, we define the random variables $U$ and $V$ of respective means $\mu_U$ and $\mu_V$, such that:
	\begin{equation}
		\begin{bmatrix}U\\V\end{bmatrix}\sim\mathcal{N}\left(\begin{bmatrix}\mu_U\\\mu_V\end{bmatrix},\begin{bmatrix}A&B\\B^\top&D\end{bmatrix}\right)
	\end{equation}
	The conditional distribution of $U$ given $V=v$ is:
	\begin{equation}
		U|V=v\sim\mathcal{N}\left(\mu_U+BD^{-1}(v-\mu_V),\:A-BD^{-1}B^\top\right)\label{eq:Gaussian_conditioning}
	\end{equation}
\end{proposition}

\begin{proof}
	Using the inferior right Schur inversion formula (see \Cref{th:Schur}) to invert the joint covariance by blocks yields:
	\begin{equation}
		\begin{bmatrix}\bar{A}&\bar{B}\\\bar{C}&\bar{D}\end{bmatrix}\triangleq\begin{bmatrix}A&B\\B^\top&D\end{bmatrix}^{-1}
	\end{equation}

	For some realisation $u$ of $U$, we define $\tilde{u}\triangleq u-\mu_U$ and $\tilde{v}\triangleq v-\mu_V$, and we compute:
	\begin{equation}
		\begin{aligned}
		\alpha&\triangleq\begin{bmatrix}\tilde{u}\\\tilde{v}\end{bmatrix}^\top\begin{bmatrix}\bar{A}&\bar{B}\\\bar{C}&\bar{D}\end{bmatrix}\begin{bmatrix}\tilde{u}\\\tilde{v}\end{bmatrix}\\
		&=\tilde{u}^\top\bar{A}\tilde{u}+\tilde{u}^\top\bar{B}\tilde{v}+\tilde{v}^\top\bar{C}\tilde{u}+\tilde{v}^\top\bar{D}\tilde{v}\\
		&=\tilde{u}^\top S^{-1}\tilde{u}+\tilde{u}^\top S^{-1}\left(-BD^{-1}\tilde{v}\right)+\left(-BD^{-1}\tilde{v}\right)^\top S^{-1}\tilde{u}+\\
		&\qquad\left(-BD^{-1}\tilde{v}\right)^\top S^{-1}\left(-BD^{-1}\tilde{v}\right)+\tilde{v}^\top D^{-1}\tilde{v}\\
		&=(\tilde{u}-BD^{-1}\tilde{v})^\top S^{-1}(\tilde{u}-BD^{-1}\tilde{v})+\tilde{v}^\top D^{-1}\tilde{v}
		\end{aligned}
	\end{equation}
	where $S\triangleq A-BD^{-1}B^\top$ is the Schur complement with respect to $D$.

	The conditional density is thus:
	\begin{equation}
		\begin{aligned}
		p\left(\tilde{u}|\tilde{v}\right)&\propto\frac{\displaystyle p\left(\tilde{u},\tilde{v}\right)}{\displaystyle p\left(\tilde{v}\right)}\\
		&\propto\frac{\displaystyle\exp\left(-\frac{1}{2}\alpha\right)}{\displaystyle\exp\left(-\frac{1}{2}\tilde{v}^\top D^{-1}\tilde{v}\right)}\\
		&\propto\exp\left(-\frac{1}{2}(\tilde{u}-BD^{-1}\tilde{v})^\top S^{-1}(\tilde{u}-BD^{-1}\tilde{v})\right)
		\end{aligned}
	\end{equation}

	With this formulation, it comes that $U|V=v$ is Gaussian with mean:
	\begin{equation}
		\mu_{U|V}\triangleq\mu_U+BD^{-1}\left(v-\mu_V\right)
	\end{equation}
	and covariance:
	\begin{equation}
		P_{U|V}\triangleq S=A-BD^{-1}B^\top
	\end{equation}
\end{proof}

\begin{remark}[Kalman gain]
	By defining the Kalman gain:
	\begin{equation}
		K\triangleq BD^{-1}
	\end{equation}
	the conditional distribution \eqref{eq:Gaussian_conditioning} can be rewritten as:
	\begin{equation}
		U|V=v\sim\mathcal{N}\left(\mu_U+K(v-\mu_V),\:A-KB^\top\right)
	\end{equation}
\end{remark}

\section{Schur Inversion Formula}

\begin{proposition}[Schur inversion formula]\label{th:Schur}
	For $n,p\in\mathbb{N}$, let $A\in\mathcal{M}_{n,n}(\mathbb{R})$, $B\in\mathcal{M}_{n,p}(\mathbb{R})$, $C\in\mathcal{M}_{p,n}(\mathbb{R})$, and $D\in\mathcal{M}_{p,p}(\mathbb{R})$. We assume that $A$ is invertible, and we define:
	\begin{equation}
		M\triangleq\begin{bmatrix}A&B\\C&D\end{bmatrix}
	\end{equation}

	Then $M$ is invertible and:
	\begin{equation}
		M^{-1}=\begin{bmatrix}A^{-1}+A^{-1}B(M/A)^{-1}CA^{-1}&-A^{-1}B(M/A)^{-1}\\-(M/A)^{-1}CA^{-1}&(M/A)^{-1}\end{bmatrix}
	\end{equation}
	where $M/A$ is the Schur complement of $M$ with respect to $A$:
	\begin{equation}
		M/A\triangleq D-CA^{-1}B
	\end{equation}
\end{proposition}

\begin{proof}
	Let $\bar{M}$ be the matrix defined by the right-hand side of the formula, and write its blocks as:
	\begin{equation}
		\begin{bmatrix}\bar{A}&\bar{B}\\\bar{C}&\bar{D}\end{bmatrix}=\bar{M}
	\end{equation}

	The identity $M\bar{M}=I$ yields:
	\begin{equation}
		\left\{\begin{aligned}
		A\bar{A}+B\bar{C}&=I\\
		A\bar{B}+B\bar{D}&=0\\
		C\bar{A}+D\bar{C}&=0\\
		C\bar{B}+D\bar{D}&=I\\
		\end{aligned}\right.
	\end{equation}

	In particular, the second equation can be rewritten as:
	\begin{equation}
		\bar{B}=-A^{-1}B\bar{D}\label{eq:bar_B}
	\end{equation}
	which can be plugged into the fourth equation to give:
	\begin{equation}
		\begin{aligned}
		&C\bar{B}+D\bar{D}=I\\
		\iff\quad&-CA^{-1}B\bar{D}+D\bar{D}=I\\
		\iff\quad&(D-CA^{-1}B)\bar{D}=I\\
		\iff\quad&(M/A)\bar{D}=I
		\end{aligned}
	\end{equation}

	Since $(M/A)$ and $\bar{D}$ are square matrices, $M/A$ is invertible, and:
	\begin{equation}
		\bar{D}=(M/A)^{-1}
	\end{equation}

	\eqref{eq:bar_B} can then be solved:
	\begin{equation}
		\bar{B}=-A^{-1}B(M/A)^{-1}
	\end{equation}

	A similar operation can be performed on the two remaining equations to obtain:
	\begin{equation}
		\bar{A}=A^{-1}(I-B\bar{C})
	\end{equation}
	and:
	\begin{equation}
		\begin{aligned}
		&C\bar{A}+D\bar{C}=0\\
		\iff\quad&CA^{-1}(I-B\bar{C})+D\bar{C}=0\\
		\iff\quad&CA^{-1}+(D-CA^{-1}B)\bar{C}=0\\
		\iff\quad&CA^{-1}+(M/A)\bar{C}=0\\
		\iff\quad&\bar{C}=-(M/A)^{-1}CA^{-1}
		\end{aligned}
	\end{equation}

	Finally:
	\begin{equation}
		\bar{A}=A^{-1}(I-B\bar{C})=A^{-1}+A^{-1}B(M/A)^{-1}CA^{-1}
	\end{equation}

	It can be verified by direct calculations that:
	\begin{equation}
		M\bar{M}=\bar{M}M=I
	\end{equation}
	which proves that $M$ is indeed invertible with inverse $\bar{M}$.
\end{proof}

\begin{remark}[Inversion with respect to $D$]
	If $D$ is assumed to be invertible instead of $A$, a block inversion formula can still be exhibited. Let $M/D$ denote the Schur complement with respect to $D$:
	\begin{equation}
		M/D\triangleq A-BD^{-1}C
	\end{equation}

	We have:
	\begin{equation}
		M^{-1}=\begin{bmatrix}(M/D)^{-1}&-(M/D)^{-1}BD^{-1}\\-D^{-1}C(M/D)^{-1}&D^{-1}+D^{-1}C(M/D)^{-1}BD^{-1}\end{bmatrix}
	\end{equation}
\end{remark}

\begin{remark}[Combining both formulations]
	When both $A$ and $D$ are invertible, the Woodbury lemma proves that the two formulations give the same matrices for each block. In particular, they can be combined block by block at will.
\end{remark}

\section{Woodbury Inversion Formula}

\begin{proposition}[Inversion identity]\label{th:simplified_Woodbury}
	For $n\in\mathbb{N}$, let $P\in\mathcal{M}_{n,n}(\mathbb{R})$ such that $I+P$ is invertible. Then:
	\begin{equation}
		(I+P)^{-1}=I-(I+P)^{-1}P
	\end{equation}
\end{proposition}

\begin{proof}
	This result is direct from:
	\begin{equation}
		I=(I+P)^{-1}(I+P)=(I+P)^{-1}+(I+P)^{-1}P
	\end{equation}
\end{proof}

\begin{proposition}[Push-through identity]\label{th:push_through}
	For $n,p\in\mathbb{N}$, let $P\in\mathcal{M}_{n,p}(\mathbb{R})$ and $Q\in\mathcal{M}_{p,n}(\mathbb{R})$ such that $I_n+PQ$ is invertible. Then $I_p+QP$ is also invertible and:
	\begin{equation}
		(I_n+PQ)^{-1}P=P(I_p+QP)^{-1}
	\end{equation}
\end{proposition}

\begin{proof}
	We start by using the identity:
	\begin{equation}
		P+PQP=P(I_p+QP)=(I_n+PQ)P
	\end{equation}

	Thus:
	\begin{equation}
		\begin{aligned}
		&&P(I_p+QP)&=(I_n+PQ)P\\
		\iff&&P&=(I_n+PQ)P(I_p+QP)^{-1}\\
		\iff&&(I_n+PQ)^{-1}P&=P(I_p+QP)^{-1}
		\end{aligned}
	\end{equation}

	This also shows that $I_p+QP$ admits a right inverse, hence is invertible.
\end{proof}

\begin{proposition}[Generalised Woodbury formula]
	For $n,p\in\mathbb{N}$, let $A\in\mathcal{M}_{n,n}(\mathbb{R})$, $B\in\mathcal{M}_{n,p}(\mathbb{R})$, $C\in\mathcal{M}_{p,n}(\mathbb{R})$, and $D\in\mathcal{M}_{p,p}(\mathbb{R})$. We assume that $A$, $D$, and $D+CA^{-1}B$ are invertible. Then $A+BD^{-1}C$ is also invertible, and:
	\begin{equation}
		\left(A+BD^{-1}C\right)^{-1}=A^{-1}-A^{-1}B\left(D+CA^{-1}B\right)^{-1}CA^{-1}
	\end{equation}
\end{proposition}

\begin{proof}
	We first factorise:
	\begin{equation}
		\begin{aligned}
		\left(A+BD^{-1}C\right)^{-1}&=\left(A(I+A^{-1}BD^{-1}C)\right)^{-1}\\
		&=\left(I+A^{-1}BD^{-1}C\right)^{-1}A^{-1}
		\end{aligned}
	\end{equation}

	Applying \Cref{th:simplified_Woodbury} to $A^{-1}BD^{-1}C$ yields:
	\begin{equation}
		\begin{aligned}
		\left(A+BD^{-1}C\right)^{-1}&=\left(I-\left(I+A^{-1}BD^{-1}C\right)^{-1}A^{-1}BD^{-1}C\right)A^{-1}\\
		&=A^{-1}-\left(I+A^{-1}BD^{-1}C\right)^{-1}A^{-1}BD^{-1}CA^{-1}
		\end{aligned}
	\end{equation}

	Using \Cref{th:push_through} with $P=A^{-1}B$ and $Q=DC$ further gives:
	\begin{equation}
		\left(I+A^{-1}BD^{-1}C\right)^{-1}A^{-1}B=A^{-1}B\left(I+D^{-1}CA^{-1}B\right)^{-1}
	\end{equation}

	Therefore:
	\begin{equation}
		\begin{aligned}
		\left(A+BD^{-1}C\right)^{-1}&=A^{-1}-A^{-1}B\left(I+D^{-1}CA^{-1}B\right)^{-1}D^{-1}CA^{-1}\\
		&=A^{-1}-A^{-1}B\left(D(I+D^{-1}CA^{-1}B)\right)^{-1}CA^{-1}\\
		&=A^{-1}-A^{-1}B\left(D+CA^{-1}B\right)^{-1}CA^{-1}
		\end{aligned}
	\end{equation}
\end{proof}

\begin{remark}[Rank reduction]
	One of the main numerical uses of the Woodbury inversion formula is to find the correct decomposition $M=A+BD^{-1}C$ such that $A$ is easy to invert, and $p$ is smaller than $n$. When these conditions are met, inverting M amounts to inverting $\tilde{M}=D+CA^{-1}B$, which has a complexity of $\mathcal{O}(p^3)$ instead of the standard $\mathcal{O}(n^3)$. This method is called inversion by rank reduction because $n$ and $p$ are the respective ranks of $M$ and $\tilde{M}$.
\end{remark}
